\chapter{Discussion and conclusion}\label{chap:conclusion}

\section{Summary of results}

The thesis set out to build a span-level hallucination detector for multilingual VLM answers that beats a judge fine-tuned on the original data alone, and to show which component brings which part of the gain. The results, by part:

\begin{itemize}
    \item \textbf{Judge backbone (Part 1).} Qwen3-VL-8B is the best of four VLM families under one fixed recipe, by more than the seed variance (H1 supported). Two fixes in the existing pipeline are worth as much as a new method. A label-aware span decoder lifts the official test score of the same generations by 0.012. A feature-coverage leak had made an earlier uncertainty pipeline look strong on validation, through free perfect scores on clean records.
    \item \textbf{Synthetic data (Part 2).} Synthetic training data does not improve the judge, even when it is realistic by a classifier two-sample test. The clean/labeled decomposition shows that its gain on hallucinated answers comes from flagging more and is paid for by false alarms on clean answers. Matching the hallucinated share removes both. H2 holds for the conditions that were run, H3(a) holds and H3(b) shows no effect.
    \item \textbf{Uncertainty and fusion (Part 3).} Four components improve the judge, each with an interpretable mechanism.
    \begin{itemize}
        \item The per-label self-consistency vote over the judge's samples improves span localisation.
        \item The CRF probe on the frozen judge is weak alone, but adds a consistent gain as one voter.
        \item Graded per-character probabilities improve the rank correlation with annotator agreement.
        \item A second probe on a frozen judge of another family (Gemma-3-12B) adds a small further gain as an eleventh voter, though not as a replacement for the first.
    \end{itemize}
    Together they reach Cor+Lbl 0.4312 on test against 0.4132 for the original-data judge (three-seed means). The gain holds for every seed and language, and on validation the final detector meets the pre-registered criterion on both halves.
    \item \textbf{Overlap of the two uncertainty signals (H4).} The voter and the probe succeed in the same category--language cells (H4 not supported at the cell level), but their record-level errors are only partly shared. A probe on a judge of another family makes errors only slightly less correlated with the judge's (0.74 against 0.76).
\end{itemize}

\section{Position in the research landscape}

The two strongest SHROOM-Visions systems converge on the same principle as our detector: agreement among several decodes becomes the character probability.
\begin{itemize}
    \item \textbf{Vroom-Vroom}~\cite{vroomvroom2026} combines five fine-tuned judges from four model families. Its committee probability, the mean of confidence times vote over judges, is the same quantity as our $\bar p_L(c)$ (Equation~\ref{eq:pbar}), and it also merges runs of equal probability into spans.
    \item \textbf{TÜRKSAT}~\cite{karamanlioglu2026turksat} uses the vote fraction of a prompted commercial model sampled 6--9 times, which corresponds to our candidate G2.
\end{itemize}
What separates their scores from ours is therefore mostly not the aggregation rule. It is the strength and diversity of the voters: larger and more varied judges, or a stronger commercial model. It is also the prompt content, including the image file name, which TÜRKSAT found to be a useful hint. Our results suggest that much of the benefit of a committee is available from \emph{one} judge, by sampling it several times and voting per label. Our four-backbone committee scored below the single-judge vote. One caveat is that three of its members used the older decoder. The top system reported that internal-activation probes did not help its committee. We find that a probe does help a self-consistency vote, by a small margin, when it is trained with multi-label targets and decodes spans with a CRF. This fits the view that probes and generative judges make correlated but not identical errors: the record-level correlation is 0.76 in our case. A probe on a judge of another family lowers it only to 0.74, yet still adds a further gain when both probes vote. Probe diversity thus helps a little, in the same way as judge diversity helps the top committee.

For the text-only predecessor task, synthetic data was a central ingredient of strong systems~\cite{rykov-etal-2025-smurfcat,mo-etal-2025-team}. Our isolated tests do not contradict this for settings with little annotated data. They show, however, that with 12k annotated multimodal answers the main effect of adding hallucinated examples is a shift of the model's prior, and that the label-dependent metric punishes this through false alarms on clean answers. A fair evaluation of synthetic data for this task must therefore control the hallucinated share and the number of training steps. Without these controls, gains on hallucinated records can look like progress while the overall score stays flat.

Methodologically, the program shows that validation reuse can be limited in practice without slowing the work much. The decision log contains nine decisions and 28 candidate evaluations. Several results that looked promising on val\_select did not confirm on val\_confirm, including the first single-label vote and, by the strict criterion, F2 and F3 alone. The test results then agreed with the order seen on validation.

\paragraph{The image file name.}
Our detector deliberately does not see the image file name, although TÜRKSAT found it useful. The reason is that the name describes how the dataset was built rather than what is in a given answer. The numbers below are from the training and validation data:
\begin{itemize}
    \item \textbf{Many names are descriptive.} 47\% of answers (45\% on test) belong to images named after the object and an unusual variant of it, for example a lion cub (\texttt{lion\_Lion\_cub}) or a peeled lemon (\texttt{lemon\_peeled\_lemon}). The builders chose these images for that property.
    \item \textbf{The name shifts the prior.} Such answers contain a mischaracterization in 54\% of cases, against 36\% for the other images.
    \item \textbf{It does not point to the span.} Among hallucinated English answers about these images, a word of the variety name falls inside a gold span in only 14\% of cases, although 46\% of the answers mention it.
\end{itemize}
A judge that reads the name therefore mainly learns that a descriptive name means a higher chance of mischaracterization. This is the same kind of construction shortcut that Section~\ref{sec:exp-synth} found in template-based synthetic data. In real use, images are rarely named by a curator, and a name that is present can be wrong. Part of the lead of the top systems may thus rest on information that does not carry over beyond this benchmark. A detector that does use such metadata should be trained with the metadata randomly removed and reported both with and without it.

\section{Comparative analysis}

Compared with our shared-task entry, the final detector gains $+0.034$ Cor+Lbl on test. About half comes from engineering of the base pipeline: the reimplementation and the decoder fix. The other half comes from the uncertainty components. On the post-competition leaderboard, our G1 submission moved the team from 8th, 6th, 7th and 5th to 5th, 5th, 5th and 2nd in English, French, Italian and Chinese; the best J1 seed would place 5th, 5th, 4th and 2nd, and the J1 three-seed mean 5th, 5th, 5th and 2nd. The remaining gap to the top teams is largest in Cor. In English, our Cor is 0.456 against 0.553 for the winner, while our IoU (0.406) is 0.08 behind its 0.487. This points to calibration of per-character probabilities, rather than span finding, as the most promising direction. The graded output already moved Cor by 0.014 without changing any span.

\section{Limitations}

\begin{itemize}
    \item \textbf{One task, one dataset.} All conclusions are drawn on SHROOM-Visions 2026. The answers were written by humans to look like VLM output~\cite{mickus2026humansdreamelectricsheep}, so the error distribution of a deployed VLM may differ.
    \item \textbf{Strength of the confirmation.} The detector meets the strict criterion against C0 only after the graded-probability step, which was decided after the main integration decision. That step and the second probe are reported as post-hoc extensions. Their candidates were registered before scoring, and their test gains have the same sign as on validation, but they were not part of the original plan. The earlier systems F2 and F3 improve on test and on val\_select, but their val\_confirm gains are not significant on their own.
    \item \textbf{Size of the last step.} The second probe meets the criterion against G1 only narrowly on val\_select, and its test gain ($+0.0016$, positive on every seed) is half its validation gain. For a deployment that cannot afford a second frozen model, G1 is the better trade-off.
    \item \textbf{Compute.} All experiments use one 24\,GB GPU. Backbones were limited to 8--12B parameters in 4-bit, Pixtral needed a reduced image resolution, and the backbone comparison used one seed per backbone.
    \item \textbf{Two probe judges, both quantised.} Only two families were tried as probe judges, and Gemma-3-12B ran in 4-bit with its features stored in 16 bits. One overflowing hidden channel had to be dropped. Full-precision features might give a stronger Gemma probe.
    \item \textbf{Scope of Part 2.} Synthetic data was tested at one dose per condition and with seed-42 screening. LLM-written injections were never trained, because the gating rule did not fire. Small effects within seed noise cannot be excluded. The negative result means ``no detectable effect at this dose and data size''.
    \item \textbf{Inference cost.} The final detector needs nine decodes of the judge and one forward pass of each of two frozen models (8B and 12B) per answer. That is more than ten times the cost of the greedy judge.
\end{itemize}

\section{Future work}

\begin{itemize}
    \item \textbf{Calibration.} Learn the mapping from votes and judge probabilities to per-character probabilities directly against annotator agreement, for example by isotonic regression on train, rather than using the mean.
    \item \textbf{More diverse probe judges.} A second family helped as an extra voter, not as a replacement. Whether more families keep adding, and whether probes from larger or full-precision judges are stronger alone, is open.
    \item \textbf{Cheaper inference.} Distil the committee of decodes into the judge, so that one decode carries the consistency signal.
    \item \textbf{Synthetic data with matched priors and long training.} Revisit synthetic data with the hallucinated share matched and real-data exposure held fixed, for example by first training on synthetic data and then on the original recipe, the one clean Part~2 design that remains untested.
\end{itemize}

\newpage
\chapter*{Acknowledgements}
\addcontentsline{toc}{chapter}{Acknowledgements}
I thank my supervisor Alexander Panchenko and co-supervisor Elisei Rykov for their guidance throughout this work and on our shared-task system. I thank the SHROOM-Visions 2026 organisers for the data, the scorer and the leaderboard, and the Skoltech NLP group for access to the GPU server on which the experiments were run.
% TODO: add funding acknowledgements if applicable.
